\documentclass[10pt]{article}
\usepackage{TLCresume}
\usepackage{microtype}
\begin{document}
\begin{tabularx}{\textwidth}{@{}p{0.24\textwidth}@{\hspace{0.03\textwidth}}X@{}}
\begin{minipage}[t]{\linewidth}
\name{Soham Datta}\\[3pt]
\meta{Pune, India}
\end{minipage}
&
\begin{minipage}[t]{\linewidth}
\raggedleft
\meta{\href{mailto:thesohamdatta@gmail.com}{thesohamdatta@gmail.com}}\\[-1pt]
\meta{+91 9420984066}\\[-1pt]
\meta{\href{https://linkedin.com/in/thesohamdatta}{linkedin.com/in/thesohamdatta}}\\[-1pt]
\meta{\href{https://github.com/thesohamdatta}{github.com/thesohamdatta}}\\[-1pt]
\meta{\href{https://sohamdatta.framer.ai}{sohamdatta.framer.ai}}
\end{minipage}
\end{tabularx}
\vspace{0.12in}

\begin{tabularx}{\textwidth}{@{}p{0.24\textwidth}@{\hspace{0.03\textwidth}}X@{}}
\sectionlabel{Experience}
&
\begin{minipage}[t]{\linewidth}
\role{Aura} \meta{\hfill Founder, Jun 2025 -- Present}\\[-1pt]
\meta{Wearable AI · Embedded Systems · Open Source}
\begin{resumeitems}
\item \body{Designed and 3D-printed a wearable pendant around Seeed Studio XIAO ESP32-S3 Sense, camera, LiPo battery, and microphone.}
\item \body{Wrote FreeRTOS firmware for 16 kHz I2S DMA audio capture, Opus compression, and OTA updates.}
\item \body{Integrated the device with the Omi Glass open-source ecosystem for audio streaming, transcription, and memory capture.}
\item \body{Diagnosed a Windows Electron/WebGL GPU memory leak in Omi (\href{https://github.com/BasedHardware/omi/issues/8438}{Issue \#8438}), resolved downstream in \href{https://github.com/BasedHardware/omi/pull/8902}{PR \#8902}. Built offline agglomerative speaker clustering with 52 unit tests (\href{https://github.com/BasedHardware/omi/pull/8919}{PR \#8919}), later adopted downstream in PR \#12471.}
\end{resumeitems}
\role{AI \& ML Intern} \meta{\hfill AICTE via Google, Oct 2024 -- Dec 2024}\\[-1pt]
\meta{Computer Vision · TensorFlow · ML Kit}
\begin{resumeitems}
\item \body{Built and evaluated lightweight computer vision models with TensorFlow and Google ML Kit for on-device classification.}
\item \body{Applied transfer learning and data augmentation under mobile memory constraints.}
\end{resumeitems}
\meta{\textbf{Past Experience}}\\[-1pt]
\body{\textbf{Associate · Reliance} · Jul 2024 -- Oct 2024. Customer service desk, complaints, memberships, vouchers, exchanges, documentation, and store-team coordination.}\\[2pt]
\body{\textbf{Research Assistant · PES Modern College} · Sep 2023 -- Mar 2024. Researched digital learning resources, organised material, and guided students through workshops and awareness activities.}
\end{minipage}
\\[-1pt]
\end{tabularx}

\begin{tabularx}{\textwidth}{@{}p{0.24\textwidth}@{\hspace{0.03\textwidth}}X@{}}
\sectionlabel{Projects}
&
\begin{minipage}[t]{\linewidth}
\role{Mia} \meta{\hfill Personal AI Engineering OS}\\[-1pt]
\meta{TypeScript · Bun · Claude · Codex · Vitest · JSONL}\\[-1pt]
\meta{\href{https://github.com/thesohamdatta/Mia}{github.com/thesohamdatta/Mia}}
\begin{resumeitems}
\item \body{Built a local-first TypeScript CLI that runs a 5-stage engineering harness: grill, spec, plan, review, ship.}
\item \body{Added unified adapters for Claude, OpenAI Codex, and local models, with lint, typecheck, Vitest gates and immutable JSONL state.}
\item \body{Replaced the background daemon architecture with stateless event-driven execution (\href{https://github.com/thesohamdatta/Mia/blob/main/docs/decisions/ADR-0001-eliminate-daemon.md}{ADR-0001}).}
\end{resumeitems}
\role{Voice AI Systems} \meta{\hfill Real-Time Agents}\\[-1pt]
\meta{LiveKit · WebRTC · Silero VAD · Python · Kotlin}\\[-1pt]
\meta{\href{https://github.com/thesohamdatta/John}{John} · \href{https://github.com/thesohamdatta/I-am-Mia}{I-am-Mia} · \href{https://github.com/thesohamdatta/voicecoder}{voicecoder}}
\begin{resumeitems}
\item \body{Built real-time conversational voice agents with LiveKit and client-side Silero VAD for barge-in and turn detection.}
\item \body{Built John, an Android client connected to a Python LiveKit agent with tool calling; built voicecoder, a VS Code extension with multi-provider fallback.}
\end{resumeitems}
\role{LLM-Council} \meta{\hfill Multi-Model Consensus}\\[-1pt]
\meta{Python · FastAPI · React · Vite · \href{https://github.com/thesohamdatta/LLM-Council}{github.com/thesohamdatta/LLM-Council}}\\[-1pt]
\body{Concurrent multi-LLM pipeline across OpenAI, Anthropic, Gemini, and Groq with Analyst, Skeptic, and Synthesizer stages.}
\end{minipage}
\\[-1pt]
\end{tabularx}

\begin{tabularx}{\textwidth}{@{}p{0.24\textwidth}@{\hspace{0.03\textwidth}}X@{}}
\sectionlabel{Skills}
&
\begin{minipage}[t]{\linewidth}
\body{\skilllabel{AI / LLM} · AI Agents · Tool Calling (MCP) · RAG · Memory Graphs · Multi-Agent Systems}\\[2pt]
\body{\skilllabel{Voice / Multimodal} · LiveKit · WebRTC · Silero VAD · Streaming STT/TTS · TensorFlow · ML Kit}\\[2pt]
\body{\skilllabel{Engineering} · Python · FastAPI · PyTorch · TypeScript · React · Kotlin · C/C++ · Git · Vitest}\\[2pt]
\body{\skilllabel{Hardware} · ESP32-S3 · FreeRTOS · 3D CAD · Rapid Prototyping}
\end{minipage}
\\[-1pt]
\end{tabularx}

\begin{tabularx}{\textwidth}{@{}p{0.24\textwidth}@{\hspace{0.03\textwidth}}X@{}}
\sectionlabel{Education}
&
\begin{minipage}[t]{\linewidth}
\role{Bachelor of Engineering (B.E.)} \meta{\hfill 2022 -- Jun 2026}\\[-1pt]
\body{Artificial Intelligence \& Machine Learning · University of Pune (SPPU) · Pune, India}
\end{minipage}
\\
\end{tabularx}
\end{document}
